\section{Обзор существующих систем физического моделирования}

На сегодняшний день существует ряд библиотек и систем физического моделирования, применяемых в 
робототехнических задачах. В данном разделе рассматриваются наиболее распространённые из них с акцентом на 
характеристики, релевантные для задач, сформулированных в разделе~1.1.

\subsection{Bullet Physics}

Bullet Physics~\cite{coumans2015bullet} --- библиотека с открытым исходным кодом, распространяемая под лицензией 
zlib. Она поддерживает симуляцию твёрдых тел, мягких тел и транспортных средств. Bullet использует импульсный 
решатель ограничений на основе алгоритма Projected Gauss-Seidel (PGS) и поддерживает несколько методов 
обнаружения столкновений.

Широкое использование Bullet в робототехнике обусловлено его интеграцией в симулятор 
Gazebo~\cite{koenig2004design} --- стандартный инструмент экосистемы ROS. Однако поддержка графического 
ускорителя в Bullet ограничена: исторически предлагалась экспериментальная реализация на CUDA для 
\textit{широкой фазы}, которая в дальнейшем не получила активного развития. Производительность "классической" 
реализации при числе тел свыше $10^3$ становится узким местом для задач обучения с подкреплением.

\subsection{NVIDIA PhysX}

NVIDIA PhysX~\cite{physx2023} --- коммерческая библиотека физического моделирования, исходный код которой 
доступен на GitHub под ограниченной лицензией. PhysX поддерживает аппаратное ускорение через CUDA, включая 
параллельную обработку тел и решение ограничений. Библиотека используется в симуляционных платформах Omniverse 
Isaac Sim, что делает её стандартом для задач робототехники с аппаратным ускорением на графическом процессоре.

Ограничения PhysX для разработки в рамках «Модейлера»: лицензия требует согласования для коммерческого и 
государственного применения; жёсткая привязка к стеку NVIDIA CUDA ограничивает поддержку оборудования; 
архитектура не рассчитана на замену отдельных компонентов собственными реализациями.

\subsection{MuJoCo}

MuJoCo (Multi-Joint dynamics with Contact)~\cite{todorov2012mujoco} --- физический движок, разработанный 
специально для задач робототехники и биомеханики. С 2021 года он доступен бесплатно; в 2022 году Google 
DeepMind открыла его исходный код.

Ключевые особенности MuJoCo:

\begin{itemize}
  \item использование обобщённых координат вместо декартовых повышает точность при моделировании кинематических цепей;
  \item собственный итерационный решатель ограничений, оптимизированный для сочленённых тел;
  \item поддержка симуляции сухого трения и деформируемых тел.
\end{itemize}

Поддержка графических ускорителей в MuJoCo реализована в виде экспериментального бэкенда MJX на основе 
JAX~\cite{mjx2023}, что ограничивает его применимость вне экосистемы Python/JAX.

\subsection{Open Dynamics Engine (ODE)}

ODE~\cite{smith2005ode} --- одна из первых библиотек физического моделирования с открытым исходным кодом, 
оказавшая значительное влияние на развитие отрасли. Она реализует решатель импульсов на основе 
LCP-формулировки (Linear Complementarity Problem) и поддерживает большое число видов сочленений.

На сегодняшний день ODE не получает активного развития: последний крупный релиз датируется 2017 годом. 
поддержка графических ускорителей отсутствует, производительность уступает современным решениям. Тем не менее 
ODE остаётся референсной реализацией для верификации результатов физического моделирования.

\subsection{Rapier}

Rapier~\cite{crozet2021rapier} --- современная библиотека физического моделирования, реализованная на языке 
Rust. Она обеспечивает высокую производительность и потокобезопасность благодаря модели владения памятью Rust. 
Rapier поддерживает параллельную обработку через SIMD-инструкции и многопоточность.

Поддержка вычислений на графических процессорах в Rapier в настоящее время не реализована. Библиотека 
ориентирована прежде всего на игровую индустрию и веб-разработку\footnote{Через трансляцию в WebAssembly}.

\subsection{NVIDIA Newton}

В 2024 году компания NVIDIA анонсировала физический движок Newton~\cite{nvidia2024newton}, разработанный 
совместно с Google DeepMind и Microsoft. Newton изначально ориентирован на использование графических 
ускорителей и задачи обучения роботов с подкреплением. Он построен на основе Warp (CUDA-ускоренный фреймворк 
на Python) и использует дифференцируемую физику для обучения.

На момент написания данной работы Newton находится на ранней стадии разработки и доступен с ограничениями. Его 
архитектура привязана к экосистеме Python/CUDA, графических ускорителям производства Nvidia и не предполагает 
использования в качестве встраиваемого компонента приложений на C++.

\subsection{Выводы по разделу}

Анализ существующих решений показывает, что ни одна из рассмотренных библиотек в полной мере не удовлетворяет 
совокупности требований проекта «Модейлер»:

\begin{itemize}
  \item свободная лицензия, допускающая коммерческое и государственное применение;
  \item полноценная поддержка вычислений на графических ускорителях разных производителей;
  \item C++ API для встраивания в мультиплатформенное настольное приложение;
  \item возможность расширения и адаптации под нужды конкретной платформы.
\end{itemize}

Это обосновывает разработку собственного физического движка \textit{dyphur}. В качестве ориентира при 
проектировании используются лучшие практики из рассмотренных систем:

\begin{itemize}
  \item позиционный метод решения ограничений (XPBD), зарекомендовавший себя в 
        симуляторах~\cite{macklin2016xpbd};
  \item параллельные BVH-алгоритмы из PhysX и Newton;
  \item открытая лицензия Apache~2.0 по образцу MuJoCo и Rapier, которая позволит сформировать вокруг проекта активное 
  сообщество.
\end{itemize}
